\section*{Reproducibility Statement}
All results are computed from committed result JSONs by \texttt{analysis/stats.py}; all
data figures are regenerated from the same JSONs by \texttt{figures/make\_figures.py}
(the conceptual schematic, Figure~\ref{fig:schematic}, by \texttt{figures/make\_schematic.py}).
Code, data-loading, pinned dependencies (\texttt{requirements.txt}), and a
\texttt{README} with exact run commands are available at
\repourl{}. The optimizer (Muon) is pinned to
commit \texttt{056a3c5869cf}. Every training and measurement run reported here was
performed on a single rented NVIDIA H100, booked through the Prime Intellect compute
exchange on hardware operated by Verda (DataCrunch). Across the $74$ committed runs the recorded
wall-clock totals $196{,}501$ seconds, or ${\approx}54.6$ H100-hours; that
counts the runs whose JSONs ship here and excludes exploratory and failed
attempts, which were not committed and whose time we did not log. (That $74$ is
coincidental and unrelated to the $74$ weight matrices of
Section~\ref{sec:general}.) This work received no external funding. Scope and caveats are stated in
Section~\ref{sec:intro} and Section~\ref{sec:limitations}.

\paragraph{Use of large language models.} A large language model was used as a
writing and analysis assistant: drafting and revising prose, generating
figure-plotting and statistics code, and cross-checking arithmetic against the
committed result files. It was not used to generate data, and no experimental
result, number, or citation in this paper originates from a model without being
recomputed from the committed JSONs or verified against the cited source. The
author is responsible for all claims.
